\subsection{Additive closures and their greatest core minorants}
\label{app:core-additive}

When owner closures are additive, pure-port compatibility has a closed
graph criterion and an explicit component-spine construction. Arbitrary
additive candidates also have a greatest core minorant whenever any
core minorant exists.

Let $c_r$ be loopless additive closures on finite nonempty $V$, so
$c_r(E)=\bigcup_{i\in E}c_r(\{i\})$. Write $i\to_r j$ when
$j\in c_r(\{i\})$; these are colored preorders, oriented from a
capturing input to its dependent source. Let $F_c$ be the reconstruction
\eqref{core:eq:closure-reconstruction}, initially allowing empty values.
Consider the two conditions
\begin{align}
 i\ne j,\ i\to_r j&\quad\Longrightarrow\quad
 c_s(\{j\})\subseteq c_r(\{i\})\quad\text{for every }s;
 \tag{D}\label{core:eq:additive-descent}\\
 i\ne j,\ i\to_r j,\ j\to_s i&\quad\Longrightarrow\quad r=s.
 \tag{R}\label{core:eq:additive-reciprocity}
\end{align}

\begin{theorem}[Additive closure characterization]
\label{core:thm:additive-closures}
The following are equivalent: $F_c$ is core; $(c_r)_r$ is the joint
pure-owner closure tuple of some core operation; and (D),(R) hold.
When they hold, every source has a source-linear minorant attaining
its complete canonical factor vector.
\end{theorem}
\begin{proof}
Necessity applies to arbitrary core operations. If $i\to_r j$ and
$k\in c_s(\{j\})\setminus c_r(\{i\})$, with $s\ne r$, use (FE)
at $A=\{i,k\}$ with one query $(s,\{k\},\{j\})$. Its cover is
empty and its sets lie in $c_r(A)$, so
$c_r(\{i,k\})=c_r(\{i\})$, a contradiction. The case $s=r$ is
idempotence, proving (D). Opposite differently colored arrows would
make $\{i\}$ and $\{j\}$ winning at their respective owners in the
neutral restriction to $\{i,j\}$. This contradicts the common-source
condition of \cref{core:thm:profiles}, proving (R).

Assume (D),(R). The uncolored union is transitive: a path
$i\to_r j\to_s k$ with $i\ne j$ keeps its first color by (D).
Its strongly connected components therefore form a partial order.
Every nontrivial component $C$ has a unique internal color $r_C$.
Indeed opposite arrows must share their color by (R); for three
points, composing through two opposite pairs with (D), then applying
(R), equates those colors. Incoming colors are uniform across a
target component. Every nontrivial proper ancestor $H$ of $C$ also
has $h\to_{r_H}C$ for every $h\in H$, by (D).

Construct $S_C$ from the $r_C$-fold of $C$ if $C$ is nontrivial,
or its single variable otherwise. Process proper ancestors of $C$
nearest first: if $H$ precedes $L$ in the ancestor order, process
$L$ before $H$. For a singleton $H=\{h\}$ attach $x_h$ at every
color $s$ with $h\to_s C$. For nontrivial $H$, first attach $x_h$
for each $h\in H$ at every such flag with $s\ne r_H$, then attach the $r_H$-fold of
all variables in $H$ at owner $r_H$.

Each target $j\in C$ occurs once, with exact factors
\begin{equation}\label{core:eq:additive-source-factor}
 D_{j,r}(E)=\bigwedge_{\substack{i\ne j\\i\to_r j}}[i\notin E].
\end{equation}
The starting fold supplies its internal filters, the singleton
attachments supply their flags, and each ancestor fold supplies its
required internal-owner flags. Every auxiliary occurrence is protected
too. For $h$ introduced in $H$, an internal predecessor has color
$r_H$ and is supplied by the $H$-fold, including for occurrences in
earlier foreign attachments. An external predecessor $i\to_r h$
lies in a proper ancestor of $H$, processed later. Condition (D)
gives $i\to_r C$, so that later attachment or fold supplies the
required filter. Thus every auxiliary factor is bounded by its
canonical factor \eqref{core:eq:additive-source-factor}.

Source decomposition gives $S_C\leq F_c$. Conversely, every winner
in $F_c$ has a supplying source $j$ in some $C$, and its unique
occurrence in $S_C$ has exactly the required factors. Hence
$F_c=\sum_C S_C$, proving core realizability and nonemptiness.
Pure probes recover the prescribed closures. The component word
for $j$ is its asserted source-linear minorant.
\end{proof}

The test also reads directly on conjunctive guard data: a literal
$[w_i\ne\{r\}]$ at source $j$ gives the nonloop arrow $i\to_r j$.
For common owner factors, (D),(R) decide joint realizability without
an auxiliary graph of occurrences or a supplied elimination order.

Now retain arbitrary colored preorders and let $\operatorname{Reach}$
be the reflexive transitive closure of their uncolored union. Call a
strongly connected component \emph{mixed} when its internal nonloop
arrows have at least two colors. Delete every mixed component and
every vertex reachable from it, and let $S$ be the surviving set.
On $S$ define the first-color completion
\begin{equation}\label{core:eq:additive-first-color}
 i\Rightarrow_r k\quad\Longleftrightarrow\quad
 i=k\ \text{ or }\
 \exists j\ne i\quad i\to_r j\ \text{ and }j\operatorname{Reach}k.
\end{equation}
Let $\bar c_r$ be the additive closures induced by these relations.

\begin{theorem}[Greatest minorant of an additive candidate]
\label{core:thm:additive-greatest-minorant}
If $S=\varnothing$, $F_c$ has no constant-free core minorant.
Otherwise its greatest core minorant, among all core operations
including nonseparable ones, is the reconstruction $F_{\bar c}$
on $S$, ignoring every variable outside $S$.
A full-support minorant exists exactly when no component is mixed.
\end{theorem}
\begin{proof}
The completion is the least extension of the original arrows satisfying
(D). A completed nonloop $r$-arrow has an original first $r$-edge;
following another arrow appends a reachable path and preserves that
edge, proving (D) and transitivity. Conversely, the uncolored union
of any (D)-closed family is transitive. An original first edge
$i\to_r j$, $i\ne j$, followed by a path from $j$ to $k$ therefore
forces $i\to_r k$. This proves minimality, including paths returning
to their starting point.

The completion preserves the original strongly connected components.
An internal completed nonloop arrow has its original first edge
inside that component. Every surviving nontrivial component has one
color, so opposite completed arrows have the same color. Thus (R)
holds and \cref{core:thm:additive-closures} constructs $G=F_{\bar c}$.
The set $S$ is predecessor-closed; a path ending in $S$ never leaves
it. No outside variable originally captures a source in $S$, and
completion retains every original incoming arrow there. Hence every
source contribution of $G$ satisfies the original guards, giving
$G\leq F_c$.

For greatestness let $T\leq F_c$ be any core word with nonempty
essential support $U$. Its support is predecessor-closed: if an
unused $i$ captured a used $j$, changing only $i$ to its capturing
atom on a pure source probe could not remove $T$'s observed winner,
contrary to $T\leq F_c$. On $U$, the singleton capture arrows of
$T$ satisfy (D),(R), even when its full closures are nonadditive.
Their strongly connected components are therefore monochromatic.
If $U$ met an original mixed component, predecessor closure would
put that entire component inside one monochromatic component of $T$,
contradicting its original colors. Nor can $U$ meet any descendant
of such a component. Thus $U\subseteq S$.

Every first-color path ending at a used source lies in $U$. By the
minimality argument for \eqref{core:eq:additive-first-color}, (D)
forces all completed incoming arrows at that source in $T$.
Its pure-owner coefficients, and therefore every actual occurrence
factor, are bounded by those of $G$. Source decomposition gives
$T\leq G$. If $S$ is empty, the nonempty support $U\subseteq S$
is impossible. If no component is mixed, $S=V$ and $G$ has full
support; a mixed component forbids full support by the same argument.
\end{proof}

For example, on $\{a,b,c\}$ take $a\leftrightarrow_0 b$ and
$a\leftrightarrow_1 c$, with other arrows reflexive. The two
monochromatic pairs form one mixed component, so there is no core
minorant. This also shows why internal component colors matter:
all its simple directed cycles are monochromatic.

The same example closes a possible shortcut in
\cref{core:thm:conditional-trace-minorants}. Its two foreign closures
pass the hereditary profile test. On the whole set their minimal
spanning coalitions are $\{a,c\},\{b,c\}$ at owner $0$ and
$\{a,b\},\{b,c\}$ at owner $1$, and every pair intersects.
On each proper retained set one closure
is identity, so the test again holds. By
\cref{core:thm:two-color-feedback} there is no raw two-color feedback,
yet the required full-support trace minorant cannot exist: its pure
source factors would be bounded by the prescribed factors, and source
decomposition would give a core minorant of this reconstruction. Thus that
trace-minorant condition is strictly stronger than absence of raw
feedback.
